% concatenated from main_arxiv.tex
\documentclass{article}

% For anonymized NeurIPS submission.
\usepackage[preprint]{neurips_2026}

\usepackage[utf8]{inputenc}
\usepackage[T1]{fontenc}
\usepackage{microtype}
\usepackage{url}
\usepackage{booktabs}
\usepackage{amsmath,amssymb,amsfonts,amsthm,mathtools}
\usepackage{graphicx}
\usepackage{xcolor}
\usepackage{hyperref}
%\usepackage[numbers,sort&compress]{natbib}

\title{Spectral Graph Sparsification Preserves Representation Geometry in Graph Neural Networks}

\author{%
  Sanjukta Krishnagopal \\
  Department of Computer Science, University of California Santa Barbara \\
  \texttt{sanjukta@ucsb.edu}
}

\newtheorem{theorem}{Theorem}
\newtheorem{proposition}{Proposition}
\newtheorem{lemma}{Lemma}
\newtheorem{corollary}{Corollary}
\newtheorem{remark}{Remark}

\begin{document}

\maketitle

\begin{abstract}
Spectral graph sparsification is a classical tool for reducing graph complexity while preserving Laplacian quadratic forms. In graph neural networks (GNNs), sparsification is often used to accelerate computation while maintaining predictive performance. In this work, we study a complementary representation-level question: does sparsification preserve the geometry of learned embeddings?

For polynomial-filter GNNs, we prove that any $\epsilon$-spectral sparsifier induces $O(\epsilon)$ perturbations in polynomial graph filters, multilayer hidden representations, and their Gram matrices. These guarantees imply stability of squared pairwise distances, class means, and covariance structure in embedding space. We further establish finite-time training stability: under smoothness and boundedness assumptions, gradient descent on dense and sparsified graphs produces weight trajectories whose separation grows at most proportionally to the sparsification distortion.

%We further establish finite-time stability of gradient descent trajectories for a one-layer model, with a multilayer extension under stronger smoothness and boundedness assumptions.

Empirically, effective-resistance sparsification validates the predicted perturbation chain on synthetic graphs and preserves hidden representation geometry on real datasets. In our experiments, the Gram matrix and training dynamics show low divergence even under substantial sparsification, consistent with the predicted stability under spectral sparsification. Hidden Gram preservation strongly predicts neighborhood preservation and class-centroid stability across FashionMNIST, Cora, and Paul15. Together, these results show that spectral sparsification preserves not only graph operators, but also the representation geometry that supports downstream use of GNN embeddings for interpretability.
\end{abstract}



\section{Introduction}
Graph neural networks (GNNs) are a standard framework for learning from relational data, with applications in citation networks, molecules, recommendation systems, and scientific data analysis
\citep{kipf2017gcn,hamilton2017graphsage,defferrard2016cnn}. Their effectiveness, however, comes with a scalability cost: training and inference often depend strongly on the size and density of the underlying graph. This is especially true for architectures whose propagation rules are expressed through graph operators, such as spectral and polynomial-filter GNNs~\cite{bruna2014spectral,defferrard2016cnn}.

Spectral graph sparsification offers a principled way to reduce this cost. Classical sparsification theory shows that a graph can be replaced by a much sparser graph whose Laplacian $\widetilde L$ satisfies
\[
(1-\varepsilon)L \preceq \widetilde L \preceq (1+\varepsilon)L,
\]
thereby preserving quadratic forms and spectral structure up to controlled distortion
\citep{spielman2011graph}. Recent work has begun to use such ideas for efficient GNNs, including effective-resistance- and Laplacian-based sparsification methods for spectral and attention-style architectures
\citep{ding2025largescale}. Most of this literature evaluates sparsification through operator approximation, runtime, or predictive accuracy.

In this paper, we ask a representation-centered question: \emph{does spectral sparsification preserve the geometry of learned hidden representations?} This question is important because GNNs are not only predictors; they are also representation learners
\citep{hamilton2017representation}. Their hidden states encode pairwise similarities, class separation, and global latent geometry. Two sparse graphs may yield similar accuracy while inducing substantially different embeddings, which can affect downstream uses such as retrieval, clustering, transfer learning, and interpretation. Recent work on feature evolution and neural collapse in GNNs further emphasizes that graph topology can shape representation geometry during training
\citep{kothapalli2023neuralcollapsegnn,papyan2020neuralcollapse}.
We focus on polynomial-filter GNNs, a broad class of spectral graph neural networks in which each layer applies a finite-degree polynomial in a graph operator before feature mixing and nonlinear activation \citep{defferrard2016cnn}. This class includes Chebyshev and other finite-hop spectral filters, and it also captures first-order propagation layers when the graph shift is written as a degree-one polynomial operator. It is therefore a natural setting for studying spectral sparsification, since the sparsifier directly perturbs the graph operator appearing in the layer.
Building on prior analyses of GNN stability and transferability under graph perturbations \citep{gama2020stability,ruiz2021gnnprocieee,ruiz2023transferability}, we move beyond operator-level guarantees and analyze hidden representation preservation directly. We prove that spectral sparsification stabilizes multilayer embeddings and their Gram matrices, yielding preservation of pairwise distances, class means, and covariance structure. We also analyze learning dynamics under gradient descent in a one-layer model and give a multilayer extension under stronger smoothness and boundedness assumptions, connecting to recent work on GNN training behavior on growing graphs \citep{krishnagopal2023gntk}.



% Graph neural networks (GNNs) are a standard framework for learning from relational data, with
% successful applications in citation networks, molecules, and recommendation systems
% \citep{kipf2017gcn,hamilton2017graphsage,defferrard2016cnn}. Their effectiveness, however, comes
% with a scalability cost: both training and inference depend strongly on the size and density of the
% underlying graph. This is particularly acute for architectures whose propagation rules are expressed
% directly in terms of graph operators, such as spectral and polynomial-filter GNNs.

% Spectral graph sparsification offers one such poprincipled way to reduce this cost. The classical sparsification theory shows
% that a graph can be replaced by a much sparser graph whose Laplacian $\widetilde L$ satisfies
% \[
% (1-\varepsilon)L \preceq \widetilde L \preceq (1+\varepsilon)L,
% \]
% thereby preserving quadratic forms and spectral structure up to controlled distortion
% \citep{spielman2011graph}. Recent work has begun to exploit this idea for GNN efficiency,
% for example through effective-resistance- and Laplacian-based sparsification for spectral or
% attention-style architectures \citep{ding2025largescale}. Most of this literature evaluates
% sparsification primarily through operator approximation, runtime, and predictive accuracy.
% In this paper, we ask a representation-centered question: \emph{Does spectral sparsification preserve the geometry of the learned hidden representations?}

% %GNNs are not only predictors but also representation learners: their hidden states encode pairwise similarity, class separation, and global latent geometry. Two sparse graphs may yield similar accuracy while inducing substantially different embeddings, which can matter for downstream uses of learned representations. Recent work on feature evolution and neural collapse in GNNs further suggests that graph topology can shape representation geometry during training.

% GNNs are not only predictors but also representation learners~\citep{hamilton2017representation}: their hidden states encode pairwise similarity, class separation, and global latent geometry.
% %Their hidden states encode pairwise similarities, class separation, and global geometric structure in latent space. Two sparse graphs can yield similar accuracy while inducing very different embeddings.
% Two sparse graphs may yield similar accuracy while inducing substantially different embeddings, which can matter for downstream uses of learned representations. Recent work on feature evolution and neural collapse in
% GNNs further emphasizes that graph topology can shape representation geometry during training
% \citep{kothapalli2023neuralcollapsegnn,papyan2020neuralcollapse}.

% Many common GNN layers can be written as polynomial graph filters followed by feature mixing and nonlinearity \citep{defferrard2016convolutional}. Building on prior analyses of GNN stability and transferability under graph perturbations \citep{gama2020stability,ruiz2021gnnprocieee,ruiz2023transferability}, we move beyond operator preservation and analyze hidden representation preservation directly. We show that spectral sparsification stabilizes multilayer embeddings and their Gram matrices, yielding preservation of pairwise distances, class means, and covariance structure. We also analyze gradient descent in a one-layer model and give a multilayer extension under stronger smoothness and boundedness assumptions, relating to recent work studying training behavior of GNNs on growing graphs \citep{krishnagopal2023gntk}.

% Empirically, synthetic sparsification sweeps validate the predicted perturbation chain, while on real datasets hidden Gram preservation strongly predicts neighborhood preservation in representation space. In the experiments, we use effective-resistance sampling as a concrete testbed, but the theory applies to any sparsification method yielding an $\varepsilon$-spectral approximation.



% Our starting point is that many common GNN layers can be written as polynomial graph filters
% followed by feature mixing and nonlinearity \citep{defferrard2016cnn}. Building on prior analyses
% of GNN stability and transferability under graph perturbations and graph shifts
% \citep{gama2020stability,ruiz2021gnnprocieee,ruiz2023transferability}, we move beyond
% operator preservation and analyze hidden representation preservation directly. We prove that
% spectral sparsification stabilizes multilayer embeddings and their Gram matrices, yielding
% preservation of pairwise distances, class means, and covariance structure. 
% We also analyze gradient descent in a one-layer model and give a multilayer extension under stronger smoothness and boundedness assumptions, showing finite-time stability of training trajectories.
% %We also analyze gradient descent in a one-layer model and show that multilayer training trajectories remain stable under sparsification-induced perturbations of the propagation operator. 
% From a broader large-graph
% perspective, our viewpoint is also related to recent work studying asymptotic and kernelized
% training behavior of GNNs on growing graphs \citep{krishnagopal2023gntk}.

% Our empirical study mirrors this theoretical picture. On synthetic graph families, filter error,
% hidden representation error, and Gram distortion scale approximately linearly with perturbation
% magnitude. On real datasets, hidden Gram preservation strongly predicts nearest-neighbor
% preservation in representation space, showing that the geometric quantities controlled by the theory
% have practical consequences for embedding-based downstream use. In the experiments, we instantiate
% the sparsifier using effective-resistance-based sampling as a concrete and principled mechanism for
% approximately obtaining spectrally similar sparse graphs; however, the theory applies to any
% sparsification method that yields an $\varepsilon$-spectral approximation.

\paragraph{Contributions.}
\begin{itemize}
    \item We introduce a representation-geometric perspective on graph sparsification for GNNs.
    \item We prove multilayer embedding stability and Gram-matrix stability under spectral sparsification.
    \item We establish finite-time training-dynamics stability in a one-layer setting, and a multilayer extension under stronger smoothness assumptions.
    \item We show empirically, using effective-resistance sparsification as a concrete testbed, on  synthetic and real data, that hidden-Gram preservation predicts neighborhood and class-geometry preservation in learned embeddings, improving interpretability and propagation fidelity.
\end{itemize}

\paragraph{Relation to prior work.}
Our work connects several lines of research. First, spectral graph sparsification provides classical
operator-level guarantees for replacing a graph by a sparse spectral approximation
\citep{spielman2011graph}. Second, recent sparsification-for-GNN methods use these ideas to
improve efficiency, focusing mainly on propagation fidelity and predictive performance
\citep{ding2025largescale}. Third, prior work studies the stability and transferability of GNNs
under graph perturbations, graph shifts, and graph-size changes
\citep{gama2020stability,ruiz2021gnnprocieee,ruiz2023transferability}. Fourth, recent work
studies the geometry of learned features in graph representation learning, including neural collapse
in GNNs \citep{kothapalli2023neuralcollapsegnn,papyan2020neuralcollapse}. Finally, our viewpoint is
also related to large-graph asymptotic analyses of GNN transferability and learning dynamics
\citep{krishnagopal2023gntk}. Our contribution lies at the intersection of these directions: we
study whether principled spectral sparsification preserves the geometry of learned representations,
not only the propagation operator or task accuracy.



\section{Preliminaries and Main Results}
\label{sec:main_results}

\subsection{Setup and Notation}

Let $G=(V,E)$ be a graph on $n$ nodes, and let $L \in \mathbb{R}^{n\times n}$ denote a symmetric positive semidefinite graph operator, such as the combinatorial or normalized Laplacian. For symmetric matrices $A,B \in \mathbb{R}^{n\times n}$, we write $A \preceq B$ if $B-A$ is positive semidefinite. We use $\|\cdot\|_2$ for the operator norm and $\|\cdot\|_F$ for the Frobenius norm.

We say that $\widetilde L$ is an $\varepsilon$-spectral sparsifier of $L$ if
\begin{equation}
(1-\varepsilon)L \preceq \widetilde L \preceq (1+\varepsilon)L,
\qquad 0<\varepsilon<1.
\label{eq:spectral_sparsifier}
\end{equation}

We study a class of $K$-layer polynomial filter graph neural networks of the form
\begin{equation}
H^{(0)} = X, \qquad
H^{(k+1)} = \sigma_k\!\bigl(p_k(L)H^{(k)}W_k\bigr),
\qquad k=0,\dots,K-1,
\label{eq:gnn_forward}
\end{equation}
where $X \in \mathbb{R}^{n\times d_0}$ is the input feature matrix, $H^{(k)} \in \mathbb{R}^{n\times d_k}$ is the hidden representation at layer $k$, $W_k \in \mathbb{R}^{d_k\times d_{k+1}}$ is the trainable weight matrix, $p_k$ is a real polynomial, and $\sigma_k$ is applied entrywise.

The corresponding network on the sparsified graph is defined by
\begin{align}
\widetilde H^{(0)} &= X, \notag\\
\widetilde H^{(k+1)}
&= \sigma_k\!\bigl(p_k(\widetilde L)\widetilde H^{(k)}W_k\bigr).
\label{eq:sparse_gnn_forward}
\end{align}
We denote the final embeddings by
\[
Z := H^{(K)},
\qquad
\widetilde Z := \widetilde H^{(K)}.
\]

Our analysis is based on the following assumptions.

\textbf{(A1) Polynomial filters.}
For each $k$, $p_k(x)=\sum_{r=0}^{d_k^\star} a_{k,r}x^r$ is a polynomial of finite degree $d_k^\star$.

\textbf{(A2) Lipschitz nonlinearities.}
Each activation $\sigma_k:\mathbb{R}\to\mathbb{R}$ is $L_{\sigma_k}$-Lipschitz.

\textbf{(A3) Bounded weights.}
There exists $B_W>0$ such that $\|W_k\|_2 \le B_W$ for all $k$.

\textbf{(A4) Bounded inputs.}
The input features satisfy $\|X\|_F \le B_X$.

\textbf{(A5) Bounded operators.}
There exists $B_L>0$ such that $\|L\|_2,\|\widetilde L\|_2 \le B_L$.

Assumptions (A1)--(A5) are standard in perturbation analyses of graph-based learning systems and are sufficient for the stability results below.

\subsection{Propagation Stability}

We first show that spectral sparsification preserves polynomial graph filters.

\begin{theorem}[Polynomial filter stability]
\label{thm:poly_stability}
Let $p(x)=\sum_{r=0}^{d} a_r x^r$ be a polynomial, and let $L,\widetilde L$ satisfy \eqref{eq:spectral_sparsifier}. Under (A5),
\begin{equation}
\|p(L)-p(\widetilde L)\|_2 \le C_p\,\varepsilon,
\label{eq:poly_stability}
\end{equation}
where
\[
C_p := \sum_{r=1}^{d} |a_r|\, r\, B_L^{\,r}.
\]
\end{theorem}

Theorem~\ref{thm:poly_stability} is the operator-level ingredient in our analysis: it says that spectral closeness of the graph implies closeness of the propagation operators used by polynomial-filter GNNs. The full proof is given in Appendix~\ref{app:poly_stability}.

\subsection{Representation Stability}

We next propagate the operator perturbation through a multilayer nonlinear architecture.

\begin{theorem}[Multilayer representation stability]
\label{thm:multilayer_stability}
Suppose that \eqref{eq:spectral_sparsifier} holds and assumptions (A1)--(A5) are satisfied. Then there exists a constant $C_{\mathrm{rep}}>0$, depending only on the network depth, filter coefficients, Lipschitz constants, and bounds $B_X,B_W,B_L$, such that
\begin{equation}
\|Z-\widetilde Z\|_F \le C_{\mathrm{rep}}\,\varepsilon.
\label{eq:multilayer_stability}
\end{equation}
\end{theorem}

%The proof is by induction over layers and uses Theorem~\ref{thm:poly_stability} together with the Lipschitz continuity of each layer map. 
The result shows that spectral sparsification controls not only the graph operator but also the learned embedding itself. The full proof is given in Appendix~\ref{app:multilayer_stability}.

\subsection{Representation Geometry}

The main conceptual contribution of this paper is that sparsification preserves the geometry of the learned embedding space.
\begin{theorem}[Gram matrix stability]
\label{thm:gram_stability}
Under the assumptions of Theorem~\ref{thm:multilayer_stability}, there exist constants
$C_{{\rm gram},2},C_{{\rm gram},F}>0$ such that
\[
\|ZZ^\top-\widetilde Z\widetilde Z^\top\|_2
\le C_{{\rm gram},2}\epsilon,
\qquad
\|ZZ^\top-\widetilde Z\widetilde Z^\top\|_F
\le C_{{\rm gram},F}\epsilon .
\]
\end{theorem}


Theorem~\ref{thm:gram_stability} upgrades representation stability to Gram-matrix stability in both spectral and Frobenius norms. Since the Gram matrix encodes all pairwise inner products between node embeddings, this implies preservation of distances and class-level statistics such as centroids and covariances in latent space. The proof is given in Appendix~\ref{app:gram_stability}.

The following corollaries formalize this implication.

\begin{corollary}[Pairwise distance preservation]
\label{cor:distance_preservation}
Let $z_i^\top$ and $\widetilde z_i^\top$ denote the $i$th rows of $Z$ and $\widetilde Z$, respectively. Then for all $i,j \in \{1,\dots,n\}$,
\begin{equation}
\bigl|
\|z_i-z_j\|_2^2 - \|\widetilde z_i-\widetilde z_j\|_2^2
\bigr|
\le 4 C_{\mathrm{gram}}\,\varepsilon.
\label{eq:distance_preservation}
\end{equation}
\end{corollary}

Now suppose that the nodes are partitioned into $C$ classes with index sets $S_1,\dots,S_C$, where $n_c := |S_c|$. Define the class mean and empirical covariance for class $c$ by
\begin{equation}
\mu_c := \frac{1}{n_c}\sum_{i\in S_c} z_i, \qquad
\Sigma_c := \frac{1}{n_c}\sum_{i\in S_c}(z_i-\mu_c)(z_i-\mu_c)^\top,
\label{eq:class_cov}
\end{equation}
and define $\widetilde\mu_c,\widetilde\Sigma_c$ analogously using $\widetilde Z$.

\begin{corollary}[Class-statistics stability]
\label{cor:class_stats}
Under the assumptions of Theorem~\ref{thm:multilayer_stability}, for each class $c$,
\begin{equation}
\|\mu_c-\widetilde\mu_c\|_2
\le \frac{1}{\sqrt{n_c}}\|Z-\widetilde Z\|_F,
\label{eq:mean_stability}
\end{equation}
and, if the row norms of $Z$ and $\widetilde Z$ are uniformly bounded by $B_Z$,
\begin{equation}
\|\Sigma_c-\widetilde\Sigma_c\|_2
\le \frac{8B_Z}{\sqrt{n_c}}\|Z-\widetilde Z\|_F.
\label{eq:cov_stability}
\end{equation}
In particular, both quantities are $O(\varepsilon)$.
\end{corollary}

Corollary~\ref{cor:class_stats} controls several representation diagnostics used in practice. Examples include the displacement of class centroids, pairwise distances and angles between centroids, within-class variance $\operatorname{tr}(\Sigma_c)$, leading covariance eigenvalues, between-/within-class scatter ratios, and nearest-centroid margins. These quantities measure whether classes remain separated, compact, and similarly oriented after sparsification. Neural-collapse-style diagnostics are also covered, since they are built from the same objects: class means, within-class covariances, and normalized centroid geometry. On bounded sets, any such diagnostic that is locally Lipschitz as a function of
$
\{\mu_c,\Sigma_c\}_{c=1}^C
$
therefore changes by at most $O(\varepsilon)$ under spectral sparsification.


% \begin{remark}[Consequences for class-geometry diagnostics]
% \label{rem:nc_stability}
% Many empirical diagnostics of class geometry---including neural-collapse-style quantities built from class means and within-class covariance matrices---are continuous, and often locally Lipschitz, functions of
% \[
% \{\mu_c,\Sigma_c\}_{c=1}^C.
% \]
% Therefore, on bounded sets, Corollary~\ref{cor:class_stats} implies that such diagnostics vary by at most $O(\varepsilon)$ under spectral sparsification.
% \end{remark}
Proofs of Corollaries~\ref{cor:distance_preservation} and~\ref{cor:class_stats} are collected in Appendix~\ref{app:geometry_corollaries}.

Theorems~\ref{thm:poly_stability}--\ref{thm:gram_stability} establish the implication chain
\[
L \approx \widetilde L
\;\Longrightarrow\;
p(L) \approx p(\widetilde L)
\;\Longrightarrow\;
Z \approx \widetilde Z
\;\Longrightarrow\;
ZZ^\top \approx \widetilde Z\widetilde Z^\top.
\]

\subsection{Training Dynamics}
We also study how sparsification perturbs training trajectories. In a one-layer model
\[
f(W;L)=\sigma\!\bigl(p(L)XW\bigr), \qquad
J(W;L)=\tfrac12\|f(W;L)-Y\|_F^2,
\]

In addition to (A1)--(A5), we make the following standard assumptions:

\textbf{(A6) Smooth activation.}
$\sigma \in C^1(\mathbb{R})$, $\sigma'$ is bounded, and $\sigma'$ is Lipschitz.

\textbf{(A7) Bounded iterates.}
The gradient descent iterates remain in a bounded set:
\[
\|W_t\|_2,\|\widetilde W_t\|_2 \le R
\qquad\text{for all } t\le T.
\]
%assume in addition to (A1)--(A5) that $\sigma\in C^1$, that $\sigma'$ is bounded and Lipschitz, and that the gradient-descent iterates remain in a bounded set up to time $T$.
\textbf[(A8)] \textbf{Bounded targets.}
The target matrix satisfies $\|Y\|_F \le B_Y$.

Let
\[
W_{t+1}=W_t-\eta \nabla_W J(W_t;L), \qquad
\widetilde W_{t+1}=\widetilde W_t-\eta \nabla_W J(\widetilde W_t;\widetilde L),
\]
with common initialization $W_0=\widetilde W_0$.

\begin{theorem}[One-layer training stability]
\label{thm:training_stability}
Under these assumptions, with constants allowed to depend on $B_Y$, there exist constants $L_W,L_S>0$ such that for all $t\le T$,
\begin{equation}
\label{eq:theorem_training_one_layer_main}
\|W_{t+1}-\widetilde W_{t+1}\|_F
\le
(1+\eta L_W)\|W_t-\widetilde W_t\|_F
+\eta L_S\|p(L)-p(\widetilde L)\|_2.
\end{equation}

Consequently, by Theorem~\ref{thm:poly_stability},
\begin{equation}
\|W_t-\widetilde W_t\|_F
\le
C_{\mathrm{train}}\varepsilon \sum_{j=0}^{t-1}(1+\eta L_W)^j,
\end{equation}
for a constant $C_{\mathrm{train}}$ independent of $\varepsilon$.
\end{theorem}
Theorem~\ref{thm:training_stability} shows that sparsification perturbs the gradient field only slightly, and hence yields finite-time stability of gradient descent trajectories in the one-layer setting. The full proof is deferred to Appendix~\ref{app:training_stability}.



\paragraph{Multilayer training dynamics.}
The one-layer result extends to full-batch multilayer gradient descent on bounded finite-time trajectories under analogous smoothness and boundedness assumptions. Let
\[
\Theta:=(W_0,\dots,W_{K-1}), \qquad
F(\Theta;L):=H^{(K)},
\]
where
\[
H^{(0)}=X, \qquad
H^{(k+1)}=\sigma_k\!\bigl(p_k(L)H^{(k)}W_k\bigr),
\]
and let
\[
J(\Theta;L):=\tfrac12\|F(\Theta;L)-Y\|_F^2.
\]
In addition to (A1)--(A5) and (A8), we make the following standard assumptions:
\begin{itemize}
    \item[(A6$'$)] \textbf{Smooth activations.}
    For each $k$, $\sigma_k \in C^1(\mathbb{R})$, $\sigma_k'$ is bounded, and $\sigma_k'$ is Lipschitz.
    \item[(A7$'$)] \textbf{Bounded iterates and hidden states.}
    There exist constants $R,B_H>0$ such that for all $t\le T$ and all $k$,
    \[
    \|W_k^{\,t}\|_2,\ \|\widetilde W_k^{\,t}\|_2 \le R,
    \qquad
    \|H^{(k),t}\|_F,\ \|\widetilde H^{(k),t}\|_F \le B_H.
    \]
\end{itemize}
Let $\Theta^t,\widetilde\Theta^t$ be the corresponding gradient descent iterates with common initialization.

\begin{proposition}[Multilayer training stability]
\label{prop:multilayer_training_stability}
Under these assumptions, there exist constants $L_\Theta,L_S>0$ such that for all $t\le T$,
\begin{equation}
\label{eq:theorem_training_multi_layer_main}
\|\Theta^{t+1}-\widetilde\Theta^{t+1}\|_F
\le
(1+\eta L_\Theta)\|\Theta^t-\widetilde\Theta^t\|_F
+
\eta L_S \sum_{k=0}^{K-1}\|p_k(L)-p_k(\widetilde L)\|_2.
\end{equation}

Consequently, by Theorem~\ref{thm:poly_stability},
\[
\|\Theta^t-\widetilde\Theta^t\|_F
\le
\eta L_S
\left(\sum_{k=0}^{K-1} C_{p_k}\right)
\varepsilon
\sum_{j=0}^{t-1}(1+\eta L_\Theta)^j
\le
C_{\mathrm{train}}^{(K)}\varepsilon
\sum_{j=0}^{t-1}(1+\eta L_\Theta)^j,
\]
where $C_{\mathrm{train}}^{(K)}$ is independent of $\varepsilon$.

\end{proposition}


The proof (see \ref{app:multilayer_training_stability}) follows the same strategy as Theorem~\ref{thm:training_stability}: one shows that the full gradient field is Lipschitz in the parameters and Lipschitz in the collection of graph filters $\{p_k(L)\}_{k=0}^{K-1}$ on bounded sets, and then combines these bounds with Theorem~\ref{thm:poly_stability}.

%\subsection{Summary}


%Thus, spectral sparsification preserves not only the graph operator used for propagation, but also the geometry of the learned representation and, in a simplified setting, the associated training dynamics.

\section{Experiments}
\label{sec:experiments}

We evaluate the central claim of the paper: in practically relevant sparsification regimes, principled graph sparsification preserves not only graph operators, but also the geometry of learned hidden representations. Our experiments combine controlled synthetic tests, which directly probe the perturbation predictions of Section~\ref{sec:main_results}, with practical representation-learning tests on real datasets.

\subsection{Setup}
\label{subsec:exp_setup_main}

We evaluate sparsification on synthetic and real graph datasets. The synthetic experiments use weighted stochastic block models (SBMs) and geometric $k$-NN graphs constructed to exhibit clear community or manifold structure. The real-data experiments use FashionMNIST \citep{xiao2017fashion}, Cora \citep{sen2008collective}, and Paul15 \citep{paul2015transcriptional}, a single-cell RNA-seq dataset profiling myeloid progenitor cells. Node features are either class-structured synthetic features, for controlled experiments, or the standard features provided with each dataset; details are given in Appendix~\ref{app:datasets}.

For the theory-aligned experiments in Figures~\ref{fig:theory_main} and~\ref{fig:training_dynamics_main}, we use deterministic GNN models to isolate the effect of graph sparsification. Specifically, we use a polynomial graph filter followed by two linear feature-mixing layers for forward-map evaluation, and one- and two-layer deterministic GNNs for training-dynamics experiments. Dense and sparsified models have matched initialization.

Sparsification is performed using effective-resistance sampling. Effective-resistance sparsification samples edges with probability proportional to their leverage scores $w_e R_e$, where $w_e$ is the edge weight and $R_e$ is the effective resistance between the endpoints of edge $e$, and then reweights sampled edges so that the sparse Laplacian remains a spectrally accurate approximation of the original Laplacian~\citep{spielman2011graph,spielman2011spectral}.
This choice provides a concrete testbed for the theory, but our results apply to any method producing a graph $\widetilde L$ satisfying the spectral comparison~\eqref{eq:spectral_sparsifier}. The sparsification level is controlled by the edge-sampling budget, with larger budgets producing denser, lower-distortion sparsifiers. Full architectural, optimization, and sparsification details are given in Appendices~\ref{app:impl} and~\ref{app:sparsifier}.

% \subsection{Evaluation metrics}
% \label{subsec:metrics_main}

% We evaluate both graph-level and representation-level distortion. At the graph level, we measure empirical spectral distortion $\epsilon_{\mathrm{emp}}$ from the generalized eigenvalue envelope of the dense and sparse Laplacians. If $\lambda_{\min}$ and $\lambda_{\max}$ are the extremal generalized eigenvalues of $(\widetilde L,L)$, we compute 
% \[
% \epsilon_{\mathrm{emp}}=\max\{1-\lambda_{\min},\lambda_{\max}-1\}.
% \]
% At the representation level, we report relative Frobenius errors of hidden embeddings $Z$ and hidden Gram matrices $ZZ^\top$. Since the Gram matrix encodes pairwise inner products, its stability directly reflects preservation of embedding geometry. For measuring finite time stability of the training dynamics, we compute the relative parameter gap given by
% \[
% \frac{\|\Theta^t-\widetilde\Theta^t\|_F}{\|\Theta^t\|_F},
% \qquad
% \|\Theta\|_F^2=\sum_k\|W_k\|_F^2,
% \]
% computed between dense-graph and sparse-graph training trajectories initialized from the same weights.

% On real datasets, we additionally measure $k$-nearest-neighbor overlap between dense and sparse embeddings, with $k=20$. 
% %This is defined as the overlap between of shared neighbors per node, averaged over held-out nodes. 
% It tests whether sparsification preserves the local geometry used by embedding-based procedures such as retrieval, local label transfer, and clustering. Finally, we also perform a knowledge-distillation experiment in Appendix~\ref{app:teacher_student}. In the knowledge-transfer task,  a dense-graph teacher GNN provides soft labels for training a student GNN on the sparsified graph. We expect transfer to succeed when sparsification preserves hidden representation geometry, because similar embeddings should induce similar class relationships, margins, and soft decision structure. Full metric definitions are given in Appendix~\ref{app:metrics}.


\subsection{Validation of representation stability}
\label{subsec:theory_validation_main}



\begin{table}[htbp!]
\centering
\setlength{\tabcolsep}{3pt}
\renewcommand{\arraystretch}{0.9}
\footnotesize
\begin{minipage}{0.47\linewidth}
\centering
\begin{tabular}{lcccccc}
\toprule
\multicolumn{7}{c}{Geometric} \\
\midrule
%target $\epsilon$         & 0.18 & 0.26 & 0.36 & 0.48 & 0.60 & 0.72 \\
$\epsilon_{\mathrm{emp}}$ & 0.37 & 0.38 & 0.41 & 0.54 & 0.66 & 0.75 \\
frac.                     & 0.90 & 0.87 & 0.87 & 0.70 & 0.58 & 0.53 \\
\bottomrule
\end{tabular}
\end{minipage}
\hfill
\begin{minipage}{0.47\linewidth}
\centering
\begin{tabular}{lcccccc}
\toprule
\multicolumn{7}{c}{SBM} \\
\midrule
%target $\epsilon$         & 0.18 & 0.26 & 0.36 & 0.48 & 0.60 & 0.72 \\
$\epsilon_{\mathrm{emp}}$ & 0.41 & 0.42 & 0.43 & 0.43 & 0.56 & 0.71 \\
frac.                     & 0.96 & 0.96 & 0.92 & 0.89 & 0.77 & 0.66 \\
\bottomrule
\end{tabular}
\end{minipage}
\caption{Mean empirical spectral distortion $\epsilon_{\mathrm{emp}}$ and retained edge fraction (frac.) for the synthetic sparsification experiments, averaged over $5$ sparsifier draws.}
\label{tab:eps_edgefrac_sidebyside}
\end{table}



\begin{figure}[htbp!]
    \centering
    \includegraphics[width=\linewidth]{Figures/fig1_theory_main.pdf}
    \caption{
\textbf{Controlled validation of representation stability.}
We measure relative polynomial-filter error, relative hidden-representation error, and relative hidden-Gram error for effective resistance sparsified weighted SBM and geometric weighted \(k\)-NN graphs. The graph operator is the scaled combinatorial Laplacian \(\bar L=L/\|L\|_2\), and the polynomial filter is \(p(\bar L)=I-0.6\bar L+0.15\bar L^2\). Across both graph families, all three quantities increase smoothly with \(\epsilon_{\mathrm{emp}}\), consistent with the stability predictions of Theorems~\ref{thm:poly_stability}--\ref{thm:gram_stability}. Markers show means over five sparsifier draws at each selected budget, and vertical error bars denote one standard deviation.
}
    \label{fig:theory_main}
\end{figure}

On weighted SBM and geometric graph families, we construct effective-resistance sparsifiers with varying degrees of sparsification, controlled by varying the sampling budget $q$. For synthetic graphs, we measure empirical spectral distortion $\epsilon_{\mathrm{emp}}$ from the generalized eigenvalue envelope of the dense and sparse Laplacians. If $\lambda_{\min}$ and $\lambda_{\max}$ are the extremal generalized eigenvalues of $(\widetilde L,L)$ on the nonconstant subspace, we compute
\[
\epsilon_{\mathrm{emp}}=\max\{1-\lambda_{\min},\lambda_{\max}-1\}.
\]
Table~\ref{tab:eps_edgefrac_sidebyside} reports the corresponding fraction of retained edges.

We then evaluate the quantities appearing in the theory: relative polynomial-filter error, relative hidden-representation error, and relative hidden-Gram error. The synthetic forward map uses the scaled Laplacian graph operator $\bar L=L/\|L\|_2$ and polynomial filter
\[
p(\bar L)=I-0.6\bar L+0.15\bar L^2.
\]
Since the Gram matrix encodes pairwise inner products, stability of $ZZ^\top$ directly reflects preservation of embedding geometry.


Figure~\ref{fig:theory_main} validates the perturbation chain of Section~\ref{sec:main_results}. Across both graph families, all three quantities increase smoothly and approximately linearly with $\epsilon_{\mathrm{emp}}$, consistent with the $O(\epsilon)$ scaling predicted by Theorems~\ref{thm:poly_stability}--\ref{thm:gram_stability}. The relative representation and Gram distortions remain small $(0.03-0.08)$ across substantial distortion regimes where only roughly half the edges remain ($\epsilon_{\mathrm{emp}} \sim 0.7 $), the relative representation and Gram distortions
remain modest. We interpret these large-$\epsilon_{\rm emp}$ points as empirical stress
tests rather than as a regime in which the perturbative constants are expected to be sharp..

\begin{figure}[htbp!]
    \centering
    \includegraphics[width=\linewidth]{Figures/fig2_training_dynamics.pdf}
\caption{
\textbf{Training-trajectory perturbation under graph sparsification.}
Relative parameter-gap trajectories between models trained on the dense graph and on effective-resistance sparsifiers for one-layer and two-layer deterministic polynomial-filter GNNs on SBM and geometric graph families. Both experiments use the same scaled polynomial-Laplacian graph shift as Figure~\ref{fig:theory_main}. Left: parameter-gap trajectories over training epochs for representative empirical distortion levels. Right: final relative parameter gap versus empirical sparsification distortion \(\epsilon_{\mathrm{emp}}\). Markers show means over five sparsifier draws at fixed target distortion, vertical error bars denote one standard deviation. The monotone increase with \(\epsilon_{\mathrm{emp}}\) is consistent with the finite-time perturbation picture of Theorem~\ref{thm:training_stability} and Proposition~\ref{prop:multilayer_training_stability}. The relative parameter gap is a trajectory-level diagnostic in parameter space and should not be interpreted as graph or representation distortion.
}
    \label{fig:training_dynamics_main}
\end{figure}

\subsection{Training dynamics under graph sparsification}
\label{subsec:training_dynamics_main}


We test whether the finite-time perturbation behavior predicted by
Theorem~\ref{thm:training_stability} and
Proposition~\ref{prop:multilayer_training_stability} is visible in controlled training experiments.
For each dense graph and sparsifier, we train dense-graph and sparse-graph models from the same initialization and compare their parameter trajectories over time. The relative parameter gap is
\[
\frac{\|\Theta^t-\widetilde\Theta^t\|_F}{\|\Theta^t\|_F},
\qquad
\|\Theta\|_F^2=\sum_k \|W_k\|_F^2,
\]
with the analogous one-layer definition when there is only one weight matrix. This is a trajectory-level diagnostic in parameter space. It should not be interpreted as graph or representation distortion.

Figure~\ref{fig:training_dynamics_main} shows that, across SBM and geometric graph families, the relative parameter gap grows smoothly over training and is consistently larger for sparsifiers with larger empirical spectral distortion \(\epsilon_{\mathrm{emp}}\). The final relative gap also increases approximately linearly with \(\epsilon_{\mathrm{emp}}\). This behavior is consistent with the recurrence in Theorem~\ref{thm:training_stability}: sparsification introduces an \(O(\epsilon)\) perturbation to the gradient field at each step, and these perturbations accumulate over finite training time.

In this corrected polynomial-Laplacian setup, the absolute parameter gaps remain small. Thus, Figure~\ref{fig:training_dynamics_main} should be read as qualitative evidence for the predicted monotone finite-time dependence on sparsification distortion, while the main representation-level validation is provided by Figure~\ref{fig:theory_main} and the real-data geometry experiments below.



% \subsection{Training dynamics under graph sparsification}
% \label{subsec:training_dynamics_main}
% We next test the stability of training dynamics in Theorem~\ref{thm:training_stability} and Proposition~\ref{prop:multilayer_training_stability}. Figure~\ref{fig:training_dynamics_main} compares parameter trajectories between models trained on the dense graph and on their effective-resistance sparsifiers, using matched initializations in both one-layer and two-layer deterministic settings on both SBM and geometric graph families.

% Across both graph families, the relative parameter gap remains small at early times and grows smoothly throughout training, with larger empirical spectral distortion $\epsilon_{\mathrm{emp}}$ producing consistently larger deviations. The final relative gap exhibits an approximately linear dependence on $\epsilon_{\mathrm{emp}}$, consistent with the finite-time perturbation bounds of Section~\ref{sec:main_results}. Notably, the two-layer model exhibits substantially smaller sensitivity to sparsification than the one-layer model, suggesting that deeper architectures may be empirically more robust to perturbations of the propagation operator.
% %We emphasize that these experiments probe a finite-time training regime rather than asymptotic convergence, and are intended as qualitative validation of the perturbation structure established in Section~2.

% %We next test whether the optimization-side stability predicted by Theorem~\ref{thm:training_stability} and Proposition~\ref{prop:multilayer_training_stability} is visible under actual sparsification. Figure~\ref{fig:training_dynamics_main} compares parameter trajectories for models trained on the dense graph and on effective-resistance sparsifiers, using matched initializations in both one-layer and two-layer settings on the SBM and geometric graph families. The relative parameter gap remains small throughout training at mild sparsification levels and grows smoothly as the empirical spectral distortion $\epsilon_{\mathrm{emp}}$ increases. Moreover, the final gap is approximately monotone in $\epsilon_{\mathrm{emp}}$, consistent with the finite-time perturbation bounds of Section~2. This complements the forward-map results of Figure~\ref{fig:theory_main} by showing that sparsification perturbs not only hidden representations but also the optimization path in a controlled way.




\subsection{Hidden-geometry preservation on real data}


\begin{figure}[htbp!]
    \centering
    \includegraphics[width=\linewidth]{Figures/fig1_final_hidden_gram_to_nn20.pdf}
    \caption{\textbf{Hidden Gram preservation predicts neighborhood preservation on real data.}
For each sparsification budget and dataset, we repeat the sparsifier with three random seeds and compute both hidden Gram distortion and mean $20$-NN overlap on held-out test nodes. Horizontal error bars show one standard deviation of hidden Gram distortion, and vertical error bars show one standard deviation of $20$-NN overlap. Across FashionMNIST, Cora, and Paul15, worse hidden Gram distortion is associated with worse neighborhood preservation, suggesting geometry preservation is a useful representation-fidelity metric.}
    \label{fig:hidden_to_nn}
\end{figure}

Our main practical question is whether preserving hidden Gram geometry also preserves useful local and class-level structure in representation space. Since the Gram matrix controls pairwise inner products, Theorem~\ref{thm:gram_stability} and Corollary~\ref{cor:distance_preservation} predict stability of pairwise embedding geometry. Corollary~\ref{cor:class_stats} further predicts stability of coarse class statistics such as centroids and within-class spread. 
For real datasets, we therefore focus directly on representation-level distortion, using
hidden Gram distortion as the geometry-preservation metric.
% For real datasets, computing the full generalized eigenvalue envelope is more expensive, so we approximate $\epsilon_{\mathrm{emp}}$ by evaluating the quadratic-form ratios
% \[
% \frac{x^\top \widetilde L x}{x^\top Lx}
% \]
% on a fixed set of random nonconstant probe vectors and taking the largest observed deviation from $1$.

Figure~\ref{fig:hidden_to_nn} tests the local-structure side of this prediction. The $20$-NN overlap measures the fraction of each node's nearest neighbors in the dense embedding that remain nearest neighbors after sparsification, averaged over held-out nodes. Across FashionMNIST, Cora, and Paul15, larger hidden Gram distortion is consistently associated with lower $20$-NN overlap. The relationship is strongly monotone on all three datasets, with Spearman correlations $-0.92$, $-0.90$, and $-0.92$, respectively. Because $20$-NN overlap measures preservation of each node's local embedding neighborhood, these results show that Gram distortion tracks a practically meaningful notion of embedding fidelity.

Figure~\ref{fig:centroids_main} tests the class-level side of the same picture. For each dataset, we fit a two-dimensional PCA projection on the dense embedding and project both dense and sparse class centroids into that same basis. Across all three datasets, sparse centroids remain close to their dense counterparts, indicating that sparsification preserves coarse class organization in addition to local neighborhoods.


Appendix~\ref{app:extra_figures} provides an additional Procrustes-aligned embedding visualization, and Appendix~\ref{app:teacher_student} evaluates another downstream diagnostic for representation geometry preservation: a knowledge-distillation task in which a dense-graph teacher provides soft labels for a student trained on the sparsified graph.







\begin{figure}[t]
    \centering
    \includegraphics[width=\linewidth]{Figures/fig3_final_centroids.pdf}
    \caption{\textbf{Class-centroid geometry under sparsification.}
    Within each dataset, both dense and sparse centroids are projected into the same PCA basis.
    Small centroid displacements indicate that sparsification preserves class geometry, consistent with the class-statistics stability in Corollary~\ref{cor:class_stats}.}
    \label{fig:centroids_main}
\end{figure}


\section{Conclusion}

We studied graph sparsification from a representation-geometric perspective. We proved that spectral sparsification preserves not only polynomial graph filters, but also multilayer hidden embeddings and their Gram matrices. These bounds imply stability of pairwise distances, class means, and covariance structure in embedding space. We also studied the training dynamics and established finite-time stability of gradient descent trajectories in controlled one-layer and multilayer settings.
Empirically, we showed that representation and training perturbations scale smoothly with empirical spectral distortion on synthetic graphs, and that hidden Gram preservation predicts neighborhood and class-geometry preservation on real datasets. Across our datasets, hidden Gram matrices, embedding-geometry diagnostics, and finite-time training trajectories exhibited substantially smaller deviations than the underlying graph-level spectral distortion. These results suggest that sparsification should be evaluated not only by efficiency or predictive accuracy, but also by the geometry it preserves. This points toward a representation-aware view of graph compression for interpretable graph learning.
\newline
\textbf{Limitations:}
Our results assume that the sparsified graph satisfies an $\epsilon$-spectral approximation, which, despite being widely used, may hold only approximately in practical implementations. The analysis also relies on boundedness, Lipschitz, and smoothness assumptions, and the training-dynamics guarantees apply to controlled finite-time full-batch settings. While we have small-scale experimental validation, larger-scale benchmarks, and more aggressive compression regimes remain important directions for future evaluations. Finally, while we show that geometric quantities are preserved and relate them to several downstream diagnostics, a deeper understanding of how representation geometry controls generalization remains an important direction for future work.
\newline
\textbf{Broader impact:}
Geometric representation preserving sparsification may reduce the computational cost of graph learning and make large-scale graph models more accessible. However, sparsification may disproportionately affect low-degree, peripheral, or minority nodes, potentially degrading performance for underrepresented groups in sensitive applications. Spectral or representation-level guarantees do not by themselves ensure fairness, privacy, or robustness.

% \section{Conclusion}

% We studied graph sparsification from a representation-geometric perspective. We prove that spectral sparsification preserves not only graph operators, but also multilayer embeddings and their Gram structure, implying stability of distances, class structure, and local neighborhoods. We also prove stability of training trajectories in one and multi-layer cases in a controlled setting.
% Empirically, we study synthetic and real datasets and show that even in moderate-distortion regimes representation and training perturbations scale smoothly with $\epsilon_{\mathrm{emp}}$, and hidden Gram preservation predicts downstream geometric fidelity on real data, including nearest neighbors. 
% These results suggest that sparsification should be evaluated not only by efficiency or accuracy, but by the geometry it preserves, pointing toward a representation-aware view of graph compression for reliable interpretation and generalization.
% \textbf{Limitations}
% Our results assume that the sparsified graph satisfies an $\epsilon$-spectral approximation, which may only hold approximately in practice. The analysis also relies on boundedness and smoothness assumptions and is evaluated in a controlled finite-time training regime, which are standard but restrictive. 
% Empirically, we focus on spectrum preserving sparsification and evaluate a limited set of architectures and datasets; behavior of alternative sparsification methods may differ, and in very aggressive sparsification, the conditions underlying the theory may degrade. Finally, while we show that geometry is preserved and study some emprical downstream tasks, a deeper understanding of how this translates to downstream generalization remains an important direction for future work.
% \textbf{Broader Impact}
% Our results support the development of more efficient graph learning systems that preserve meaningful structure, which is valuable in large-scale applications. However, sparsification may disproportionately affect low-degree or minority nodes, potentially degrading performance on underrepresented groups. While principled methods can preserve geometry, they do not guarantee uniform performance across tasks or populations.


\paragraph{Code availability.}
Code is available at \url{github.com/chimeraki/Spectral-Sparsification-Preserves-Representation-Geometry}.

%\paragraph{LLM use statement.} LLMs were used for drafting and editing, formatting, and assisting with writing code. They were not used as part of the core scientific methodology.

\bibliography{references}
\bibliographystyle{plainnat}

\appendix

\section{Proofs}
\label{app:proofs}

This appendix contains full proofs of the results stated in Section~\ref{sec:main_results}. We retain the notation introduced there.

\subsection{Auxiliary bounds}
\label{app:auxiliary}

We first record a standard identity for matrix powers.

\begin{lemma}[Telescoping identity]
\label{lem:telescoping}
For any square matrices $A,B$ and any integer $r\ge 1$,
\[
A^r - B^r = \sum_{j=0}^{r-1} A^{r-1-j}(A-B)B^j.
\]
\end{lemma}

\begin{proof}
The claim follows by induction on $r$. The case $r=1$ is immediate. If it holds for $r$, then
\[
A^{r+1}-B^{r+1}
=
A(A^r-B^r)+(A-B)B^r,
\]
and substituting the induction hypothesis yields the result.
\end{proof}

We will also use the following consequence of spectral sparsification.

\begin{lemma}
\label{lem:operator_diff}
If $(1-\varepsilon)L \preceq \widetilde L \preceq (1+\varepsilon)L$ and $\|L\|_2\le B_L$, then
\[
\|L-\widetilde L\|_2 \le \varepsilon B_L.
\]
\end{lemma}

\begin{proof}
For any unit vector $x$,
\[
-\epsilon x^\top Lx \le x^\top(\widetilde L-L)x
\le \epsilon x^\top Lx.
\]
Since $0\preceq L$ and $\|L\|_2\le B_L$, we have
$x^\top Lx\le B_L\|x\|_2^2=B_L$. Hence
\[
|x^\top(\widetilde L-L)x|\le \epsilon B_L
\]
for every unit vector $x$. Because $\widetilde L-L$ is symmetric,
\[
\|\widetilde L-L\|_2
=
\sup_{\|x\|_2=1}|x^\top(\widetilde L-L)x|
\le \epsilon B_L .
\]
\end{proof}


\subsection{Proof of Theorem~\ref{thm:poly_stability}}
\label{app:poly_stability}

\begin{proof}[Proof of Theorem~\ref{thm:poly_stability}]
Let
\[
p(x)=\sum_{r=0}^d a_r x^r.
\]
Then
\[
p(L)-p(\widetilde L)=\sum_{r=1}^d a_r (L^r-\widetilde L^r).
\]
By Lemma~\ref{lem:telescoping},
\[
L^r-\widetilde L^r
=
\sum_{j=0}^{r-1} L^{\,r-1-j}(L-\widetilde L)\widetilde L^{\,j}.
\]
Taking operator norms and using $\|L\|_2,\|\widetilde L\|_2\le B_L$,
\[
\|L^r-\widetilde L^r\|_2
\le
\sum_{j=0}^{r-1}\|L\|_2^{r-1-j}\|L-\widetilde L\|_2\|\widetilde L\|_2^j
\le
r B_L^{\,r-1}\|L-\widetilde L\|_2.
\]
Applying Lemma~\ref{lem:operator_diff},
\[
\|L^r-\widetilde L^r\|_2
\le
r B_L^{\,r}\varepsilon.
\]
Therefore,
\[
\|p(L)-p(\widetilde L)\|_2
\le
\sum_{r=1}^d |a_r|\,\|L^r-\widetilde L^r\|_2
\le
\left(\sum_{r=1}^d |a_r|\,r\,B_L^{\,r}\right)\varepsilon.
\]
This is exactly \eqref{eq:poly_stability}.
\end{proof}

\subsection{Proof of Theorem~\ref{thm:multilayer_stability}}
\label{app:multilayer_stability}

\begin{proof}[Proof of Theorem~\ref{thm:multilayer_stability}]
For each layer define
\[
A_k := p_k(L),
\qquad
\widetilde A_k := p_k(\widetilde L).
\]
Since $\|L\|_2,\|\widetilde L\|_2\le B_L$ and each $p_k$ has fixed finite degree, there exist constants
\[
M_k := \sup_{\|S\|_2\le B_L}\|p_k(S)\|_2 < \infty
\]
such that
\[
\|A_k\|_2,\|\widetilde A_k\|_2 \le M_k .
\]
We first record uniform forward bounds. Let $B_0:=B_X$. Since $\sigma_k$ is Lipschitz,
\[
\|H^{(k+1)}\|_F
\le
L_{\sigma_k}\|A_k H^{(k)}W_k\|_F
+
|\sigma_k(0)|\sqrt{n d_{k+1}} .
\]
Using $\|A_k\|_2\le M_k$ and $\|W_k\|_2\le B_W$, we get
\[
\|H^{(k+1)}\|_F
\le
L_{\sigma_k}M_k B_W \|H^{(k)}\|_F
+
|\sigma_k(0)|\sqrt{n d_{k+1}} .
\]
The same bound holds for $\widetilde H^{(k+1)}$. Hence, by induction, there are constants
$B_k'$ independent of $\epsilon$ such that
\[
\|H^{(k)}\|_F,\|\widetilde H^{(k)}\|_F\le B_k'
\qquad
\text{for all }k=0,\ldots,K .
\]
By Theorem~\ref{thm:poly_stability},
\[
\|A_k-\widetilde A_k\|_2 \le C_{p_k}\varepsilon.
\]
Let
\[
\Delta_k := \|H^{(k)}-\widetilde H^{(k)}\|_F.
\]
Since $H^{(0)}=\widetilde H^{(0)}=X$, we have $\Delta_0=0$.

Using the layer definitions, Lipschitz continuity of $\sigma_k$, and adding and subtracting $\widetilde A_k H^{(k)}W_k$, we obtain
\begin{align}
\Delta_{k+1}
&= \|\sigma_k(A_k H^{(k)}W_k)-\sigma_k(\widetilde A_k \widetilde H^{(k)}W_k)\|_F \notag\\
&\le L_{\sigma_k}\Bigl(
\|(A_k-\widetilde A_k)H^{(k)}W_k\|_F
+
\|\widetilde A_k(H^{(k)}-\widetilde H^{(k)})W_k\|_F
\Bigr).
\label{eq:delta_recursion_step}
\end{align}
By Lipschitz continuity of $\sigma_k$ and by adding and subtracting
$\widetilde A_k H^{(k)}W_k$, we have
\[
\begin{aligned}
\Delta_{k+1}
&\le
L_{\sigma_k}
\left(
\|(A_k-\widetilde A_k)H^{(k)}W_k\|_F
+
\|\widetilde A_k(H^{(k)}-\widetilde H^{(k)})W_k\|_F
\right)  \\
&\le
L_{\sigma_k}B_W
\left(
\|A_k-\widetilde A_k\|_2 \|H^{(k)}\|_F
+
\|\widetilde A_k\|_2 \Delta_k
\right).
\end{aligned}
\]
Using Theorem~\ref{thm:poly_stability} and the uniform forward bounds gives
\[
\Delta_{k+1}\le \alpha_k\epsilon+\beta_k\Delta_k ,
\]
where
\[
\alpha_k := L_{\sigma_k}B_W C_{p_k}B_k',
\qquad
\beta_k := L_{\sigma_k}B_W M_k .
\]
Since $\Delta_0=0$, iterating this finite recursion gives
\[
\Delta_K\le C_{\rm rep}\epsilon ,
\]
where $C_{\rm rep}$ depends only on the stated network and boundedness constants.
which is exactly \eqref{eq:multilayer_stability}.
\end{proof}

\subsection{Proof of Theorem~\ref{thm:gram_stability}}
\label{app:gram_stability}

\begin{proof}
Let $\Delta:=Z-\widetilde Z$. Then
\[
ZZ^\top-\widetilde Z\widetilde Z^\top
=
Z\Delta^\top+\Delta \widetilde Z^\top .
\]
For the spectral norm,
\[
\begin{aligned}
\|ZZ^\top-\widetilde Z\widetilde Z^\top\|_2
&\le
\|Z\Delta^\top\|_2+\|\Delta\widetilde Z^\top\|_2 \\
&\le
(\|Z\|_2+\|\widetilde Z\|_2)\|\Delta\|_F .
\end{aligned}
\]
For the Frobenius norm,
\[
\begin{aligned}
\|ZZ^\top-\widetilde Z\widetilde Z^\top\|_F
&\le
\|Z\Delta^\top\|_F+\|\Delta\widetilde Z^\top\|_F \\
&\le
(\|Z\|_2+\|\widetilde Z\|_2)\|\Delta\|_F .
\end{aligned}
\]
By the boundedness argument used in the proof of Theorem~\ref{thm:multilayer_stability}, both $\|Z\|_2$ and $\|\widetilde Z\|_2$ are bounded independently of $\varepsilon$. Theorem~\ref{thm:multilayer_stability} gives
$\|\Delta\|_F\le C_{\rm rep}\epsilon$. Therefore both claimed bounds follow.

\end{proof}

\subsection{Proofs of Corollaries~\ref{cor:distance_preservation} and~\ref{cor:class_stats}}
\label{app:geometry_corollaries}

\begin{proof}[Proof of Corollary~\ref{cor:distance_preservation}]
Recall that
\[
\|z_i-z_j\|_2^2
=
\langle z_i,z_i\rangle+\langle z_j,z_j\rangle-2\langle z_i,z_j\rangle.
\]
The same identity holds for $\widetilde z_i,\widetilde z_j$. Therefore,
\begin{align*}
\bigl|
\|z_i-z_j\|_2^2-\|\widetilde z_i-\widetilde z_j\|_2^2
\bigr|
&\le
|\langle z_i,z_i\rangle-\langle \widetilde z_i,\widetilde z_i\rangle|
+
|\langle z_j,z_j\rangle-\langle \widetilde z_j,\widetilde z_j\rangle| \\
&\quad +
2|\langle z_i,z_j\rangle-\langle \widetilde z_i,\widetilde z_j\rangle|.
\end{align*}
Each inner-product difference is an entry of the matrix $ZZ^\top-\widetilde Z\widetilde Z^\top$, hence is bounded by its spectral norm. Applying the spectral-norm bounds in Theorem~\ref{thm:gram_stability} gives the result.
\end{proof}

\begin{proof}[Proof of Corollary~\ref{cor:class_stats}]
Let $\Delta := Z-\widetilde Z$, and let $\Delta_i$ denote its $i$th row. Then
\[
\mu_c-\widetilde\mu_c
=
\frac{1}{n_c}\sum_{i\in S_c}\Delta_i.
\]
By Cauchy--Schwarz,
\[
\|\mu_c-\widetilde\mu_c\|_2
\le
\frac{1}{n_c}\sum_{i\in S_c}\|\Delta_i\|_2
\le
\frac{1}{\sqrt{n_c}}
\left(\sum_{i\in S_c}\|\Delta_i\|_2^2\right)^{1/2}
\le
\frac{1}{\sqrt{n_c}}\|\Delta\|_F.
\]
This proves \eqref{eq:mean_stability}.

For the covariance bound, define
\[
a_i := z_i-\mu_c,
\qquad
\widetilde a_i := \widetilde z_i-\widetilde\mu_c.
\]
Then
\[
\Sigma_c-\widetilde\Sigma_c
=
\frac{1}{n_c}\sum_{i\in S_c}(a_i a_i^\top-\widetilde a_i \widetilde a_i^\top).
\]
Using
\[
a_i a_i^\top-\widetilde a_i \widetilde a_i^\top
=
a_i(a_i-\widetilde a_i)^\top + (a_i-\widetilde a_i)\widetilde a_i^\top,
\]
we obtain
\[
\|a_i a_i^\top-\widetilde a_i \widetilde a_i^\top\|_2
\le
(\|a_i\|_2+\|\widetilde a_i\|_2)\|a_i-\widetilde a_i\|_2.
\]
If all row norms of $Z$ and $\widetilde Z$ are bounded by $B_Z$, then
\[
\|a_i\|_2,\|\widetilde a_i\|_2 \le 2B_Z.
\]
Moreover,
\[
\|a_i-\widetilde a_i\|_2
\le
\|z_i-\widetilde z_i\|_2+\|\mu_c-\widetilde\mu_c\|_2.
\]
Therefore,
\[
\|\Sigma_c-\widetilde\Sigma_c\|_2
\le
\frac{4B_Z}{n_c}\sum_{i\in S_c}
\left(\|z_i-\widetilde z_i\|_2+\|\mu_c-\widetilde\mu_c\|_2\right).
\]
Using Cauchy--Schwarz for the first term and the mean bound for the second term,
\[
\frac{1}{n_c}\sum_{i\in S_c}\|z_i-\widetilde z_i\|_2
\le
\frac{1}{\sqrt{n_c}}\|\Delta\|_F,
\qquad
\|\mu_c-\widetilde\mu_c\|_2
\le
\frac{1}{\sqrt{n_c}}\|\Delta\|_F.
\]
Combining these yields
\[
\|\Sigma_c-\widetilde\Sigma_c\|_2
\le
\frac{8B_Z}{\sqrt{n_c}}\|\Delta\|_F.
\]
Finally, Theorem~\ref{thm:multilayer_stability} gives $\|\Delta\|_F=O(\varepsilon)$.
\end{proof}

\subsection{Proof of Theorem~\ref{thm:training_stability}}
\label{app:training_stability}

We first establish two Lipschitz properties of the gradient.

\begin{lemma}[Gradient Lipschitz in weights]
\label{lem:grad_lipschitz_W}
Under (A1), (A5), (A6), and (A7), there exists $L_W>0$ such that
\[
\|\nabla_W J(W;L)-\nabla_W J(V;L)\|_F
\le
L_W \|W-V\|_F
\]
for all $W,V$ in the bounded region $\{\|W\|_2,\|V\|_2\le R\}$.
\end{lemma}

\begin{proof}
Let $S:=p(L)$. The gradient can be written as
\[
\nabla_W J(W;L)
=
X^\top S\Bigl((\sigma(SXW)-Y)\odot \sigma'(SXW)\Bigr).
\]
This expression is a composition of affine maps, matrix products, and pointwise $C^1$ nonlinearities. The Lipschitz assumption on $\sigma_k'$ ensures that these compositions have uniformly
bounded first derivatives on the prescribed bounded parameter/state region.
Under assumptions (A4)--(A8), all quantities appearing in the gradient remain in a bounded set, and on bounded sets both $\sigma$ and $\sigma'$ are Lipschitz. Hence the map $W\mapsto \nabla_W J(W;L)$ is locally Lipschitz on that bounded set. Therefore there exists $L_W>0$ such that
\[
\|\nabla_W J(W;L)-\nabla_W J(V;L)\|_F
\le
L_W\|W-V\|_F.
\]
\end{proof}

\begin{lemma}[Gradient Lipschitz in the graph operator]
\label{lem:grad_lipschitz_S}
Under (A1), (A4)--(A7), there exists $L_S>0$ such that
\[
\|\nabla_W J(W;L)-\nabla_W J(W;\widetilde L)\|_F
\le
L_S \|p(L)-p(\widetilde L)\|_2
\]
for all $W$ with $\|W\|_2\le R$.
\end{lemma}

\begin{proof}
Writing $S:=p(L)$ and $\widetilde S:=p(\widetilde L)$, the same gradient formula shows that $\nabla_W J(W;\cdot)$ is locally Lipschitz in the propagation operator on bounded sets. Since $W$ remains bounded and $\sigma,\sigma'$ are bounded and Lipschitz on bounded sets, there exists $L_S>0$ such that
\[
\|\nabla_W J(W;L)-\nabla_W J(W;\widetilde L)\|_F
\le
L_S\|p(L)-p(\widetilde L)\|_2.
\]
\end{proof}

\begin{proof}[Proof of Theorem~\ref{thm:training_stability}]
Starting from the gradient descent updates,
\[
W_{t+1}-\widetilde W_{t+1}
=
W_t-\widetilde W_t
-
\eta\Bigl(
\nabla_W J(W_t;L)-\nabla_W J(\widetilde W_t;\widetilde L)
\Bigr).
\]
Taking Frobenius norms and using the triangle inequality,
\[
\|W_{t+1}-\widetilde W_{t+1}\|_F
\le
\|W_t-\widetilde W_t\|_F
+
\eta
\|\nabla_W J(W_t;L)-\nabla_W J(\widetilde W_t;\widetilde L)\|_F.
\]
Split the gradient difference as
\begin{align*}
\nabla_W J(W_t;L)-\nabla_W J(\widetilde W_t;\widetilde L)
&=
\bigl(
\nabla_W J(W_t;L)-\nabla_W J(\widetilde W_t;L)
\bigr) \\
&\quad +
\bigl(
\nabla_W J(\widetilde W_t;L)-\nabla_W J(\widetilde W_t;\widetilde L)
\bigr).
\end{align*}
Applying Lemmas~\ref{lem:grad_lipschitz_W} and \ref{lem:grad_lipschitz_S},
\[
\|W_{t+1}-\widetilde W_{t+1}\|_F
\le
(1+\eta L_W)\|W_t-\widetilde W_t\|_F
+
\eta L_S \|p(L)-p(\widetilde L)\|_2.
\]
This proves \eqref{eq:theorem_training_one_layer_main}. Since $W_0=\widetilde W_0$, unrolling the recursion yields
\[
\|W_t-\widetilde W_t\|_F
\le
\eta L_S \|p(L)-p(\widetilde L)\|_2
\sum_{j=0}^{t-1}(1+\eta L_W)^j.
\]
Finally, Theorem~\ref{thm:poly_stability} implies
\[
\|p(L)-p(\widetilde L)\|_2 \le C_p \varepsilon,
\]
which gives the claimed bound.
\end{proof}

\subsection{Proof of Proposition~\ref{prop:multilayer_training_stability}}
\label{app:multilayer_training_stability}

In this subsection we prove the multilayer analogue of Theorem~\ref{thm:training_stability}. For a parameter tuple
\[
\Theta=(W_0,\dots,W_{K-1}),
\]
define the hidden states recursively by
\begin{align}
H^{(0)}&=X, \notag\\
H^{(k+1)}&=\sigma_k\!\bigl(S_k H^{(k)}W_k\bigr),
\qquad
S_k:=p_k(L),
\label{eq:app_multilayer_forward_dense}
\end{align}
and similarly, on the sparsified graph,
\begin{align}
\widetilde H^{(0)}&=X, \notag\\
\widetilde H^{(k+1)}
&=
\sigma_k\!\bigl(\widetilde S_k \widetilde H^{(k)}W_k\bigr),
\qquad
\widetilde S_k:=p_k(\widetilde L).
\label{eq:app_multilayer_forward_sparse}
\end{align}
Let
\[
F(\Theta;L):=H^{(K)},
\qquad
J(\Theta;L):=\frac12 \|F(\Theta;L)-Y\|_F^2,
\]
and equip parameter space with the product Frobenius norm
\[
\|\Theta\|_F^2:=\sum_{k=0}^{K-1}\|W_k\|_F^2.
\]

For each layer $k$, define the preactivation
\[
U_k:=S_k H^{(k)}W_k.
\]
We also introduce the backpropagated state gradients
\[
\Delta_K := H^{(K)}-Y,
\]
and for $k=K-1,\dots,0$,
\begin{equation}
\Delta_k
=
S_k^\top
\Bigl(
\bigl(\Delta_{k+1}\odot \sigma_k'(U_k)\bigr)W_k^\top
\Bigr).
\label{eq:app_delta_recursion}
\end{equation}
With this notation, the gradient with respect to $W_k$ is
\begin{equation}
\nabla_{W_k}J(\Theta;L)
=
(H^{(k)})^\top S_k^\top
\bigl(\Delta_{k+1}\odot \sigma_k'(U_k)\bigr).
\label{eq:app_multilayer_weight_gradient}
\end{equation}

We first prove two Lipschitz properties of the full gradient.

\begin{lemma}[Gradient Lipschitz in the parameters]
\label{lem:multilayer_grad_lipschitz_theta}
Under assumptions (A1)--(A5), (A6$'$), (A7$'$), and (A8) there exists a constant $L_\Theta>0$ such that
\[
\|\nabla_\Theta J(\Theta;L)-\nabla_\Theta J(\Phi;L)\|_F
\le
L_\Theta \|\Theta-\Phi\|_F
\]
for all parameter tuples $\Theta,\Phi$ in the bounded region determined by (A7$'$).
\end{lemma}

\begin{proof}
From the forward recursion \eqref{eq:app_multilayer_forward_dense}, the backward recursion \eqref{eq:app_delta_recursion}, and the gradient formula \eqref{eq:app_multilayer_weight_gradient}, all forward states, backward states, and gradient blocks are compositions of affine maps, matrix products, and pointwise $C^1$ nonlinearities. Under assumptions (A1)--(A5), (A6$'$), (A7$'$), and (A8) these quantities remain in a bounded set, on which all such maps are Lipschitz. Therefore the full gradient map $\Theta\mapsto \nabla_\Theta J(\Theta;L)$ is locally Lipschitz on that bounded set, so there exists $L_\Theta>0$ such that
\[
\|\nabla_\Theta J(\Theta;L)-\nabla_\Theta J(\Phi;L)\|_F
\le
L_\Theta \|\Theta-\Phi\|_F.
\]
\end{proof}

\begin{lemma}[Gradient Lipschitz in the graph filters]
\label{lem:multilayer_grad_lipschitz_filters}
Under assumptions (A1)--(A5), (A6$'$), and (A7$'$), there exists a constant $L_S>0$ such that
\[
\|\nabla_\Theta J(\Theta;L)-\nabla_\Theta J(\Theta;\widetilde L)\|_F
\le
L_S \sum_{k=0}^{K-1}\|S_k-\widetilde S_k\|_2
\]
for all parameter tuples $\Theta$ in the bounded region determined by (A7$'$).
\end{lemma}

\begin{proof}
Fix $\Theta$. Using the same parameter tuple on both graphs, the only perturbation comes from replacing $\{S_k\}_{k=0}^{K-1}$ by $\{\widetilde S_k\}_{k=0}^{K-1}$. The forward states, backward states, and gradient blocks are again compositions of affine maps, matrix products, and pointwise $C^1$ nonlinearities, and remain bounded under (A1)--(A5), (A6$'$), and (A7$'$). Hence the full gradient is locally Lipschitz in the collection of graph filters on bounded sets. Therefore there exists $L_S>0$ such that
\[
\|\nabla_\Theta J(\Theta;L)-\nabla_\Theta J(\Theta;\widetilde L)\|_F
\le
L_S \sum_{k=0}^{K-1}\|S_k-\widetilde S_k\|_2.
\]
\end{proof}

\begin{proof}[Proof of Proposition~\ref{prop:multilayer_training_stability}]
This is the same perturbation argument as in Theorem~\ref{thm:training_stability}, now applied in product parameter space and with the graph perturbation measured through the collection $\{p_k(L)-p_k(\widetilde L)\}_{k=0}^{K-1}$.

Starting from the gradient descent updates,
\begin{equation}
\Theta^{t+1}-\widetilde\Theta^{t+1}
=
\Theta^t-\widetilde\Theta^t
-
\eta\Bigl(
\nabla_\Theta J(\Theta^t;L)
-
\nabla_\Theta J(\widetilde\Theta^t;\widetilde L)
\Bigr).
\label{eq:app_multilayer_gd_difference}
\end{equation}
Taking Frobenius norms and using the triangle inequality,
\begin{equation}
\|\Theta^{t+1}-\widetilde\Theta^{t+1}\|_F
\le
\|\Theta^t-\widetilde\Theta^t\|_F
+
\eta
\|\nabla_\Theta J(\Theta^t;L)
-
\nabla_\Theta J(\widetilde\Theta^t;\widetilde L)\|_F.
\label{eq:app_multilayer_basic_recursion}
\end{equation}
Split the gradient difference as
\begin{align}
\nabla_\Theta J(\Theta^t;L)-\nabla_\Theta J(\widetilde\Theta^t;\widetilde L)
&=
\Bigl(
\nabla_\Theta J(\Theta^t;L)
-
\nabla_\Theta J(\widetilde\Theta^t;L)
\Bigr) \notag\\
&\quad+
\Bigl(
\nabla_\Theta J(\widetilde\Theta^t;L)
-
\nabla_\Theta J(\widetilde\Theta^t;\widetilde L)
\Bigr).
\label{eq:app_multilayer_gradient_split}
\end{align}
Applying Lemmas~\ref{lem:multilayer_grad_lipschitz_theta} and~\ref{lem:multilayer_grad_lipschitz_filters} gives
\[
\|\Theta^{t+1}-\widetilde\Theta^{t+1}\|_F
\le
(1+\eta L_\Theta)\|\Theta^t-\widetilde\Theta^t\|_F
+
\eta L_S
\sum_{k=0}^{K-1}\|S_k-\widetilde S_k\|_2.
\]
This proves \eqref{eq:theorem_training_multi_layer_main}.

Since $\Theta^0=\widetilde\Theta^0$, unrolling the recursion yields
\begin{equation}
\label{eq:theta}
\|\Theta^t-\widetilde\Theta^t\|_F
\le
\eta L_S
\left(
\sum_{k=0}^{K-1}\|S_k-\widetilde S_k\|_2
\right)
\sum_{j=0}^{t-1}(1+\eta L_\Theta)^j.
\end{equation}
Finally, by Theorem~\ref{thm:poly_stability}, for each $k$,
\[
\|S_k-\widetilde S_k\|_2
=
\|p_k(L)-p_k(\widetilde L)\|_2
\le
C_{p_k}\varepsilon.
\]
Therefore
\[
\sum_{k=0}^{K-1}\|S_k-\widetilde S_k\|_2
\le
\left(\sum_{k=0}^{K-1}C_{p_k}\right)\varepsilon,
\]
Substituting this estimate into Eqn.~\ref{eq:theta} gives
\[
\|\Theta^t-\widetilde\Theta^t\|_F
\le
\eta L_S
\left(\sum_{k=0}^{K-1} C_{p_k}\right)
\varepsilon
\sum_{j=0}^{t-1}(1+\eta L_\Theta)^j,
\]
which proves the claimed bound with
\[
C_{\mathrm{train}}^{(K)}
=
\eta L_S\sum_{k=0}^{K-1} C_{p_k}.
\]
\end{proof}

\section{Additional Experimental Details}
\label{app:experiments}

\subsection{Dataset construction}
\label{app:datasets}

\paragraph{Synthetic SBM.}
We generate a four-block weighted stochastic block model with block sizes \((80,80,80,80)\), within-block edge probability \(0.20\), and cross-block edge probability \(0.05\).

\paragraph{Synthetic geometric graph.}
We generate four Gaussian feature clusters with \(80\) nodes per class and \(20\)-dimensional features. A weighted \(k\)-NN graph with \(k=30\) is built from these features using Gaussian kernel weights.

\paragraph{FashionMNIST.}
We load the training split, subsample $220$ examples per class in the full setting, flatten the images, reduce to $50$ PCA dimensions, and build a weighted $15$-NN graph.

\paragraph{Cora.}
We use the Planetoid Cora benchmark and take the citation graph as the adjacency.

\paragraph{Paul15.}
We use the Paul15 single-cell dataset from Scanpy. Labels are extracted from an available cluster annotation, features are reduced to at most $64$ dimensions when needed, and a weighted $15$-NN graph is built from expression features.

\paragraph{Splits.}
%\textcolor{red}{[What does this mean? Is this a standard CS thing]}
Class-balanced random masks are used throughout. For FashionMNIST, we use $20$ train and $40$ validation examples per class. For Cora, we use $20$ train and $50$ validation examples per class. For Paul15, we use $20$ train and $30$ validation examples per class. Remaining nodes are used for testing. On synthetic datasets, analogous class-balanced masks are used.

\subsection{Architectures and optimization}
\label{app:impl}
\paragraph{Synthetic theory-facing models.}
For the controlled sparsification experiments of Figure~\ref{fig:theory_main}, we use a fixed deterministic two-stage polynomial-filter map so that the forward distortion induced by sparsification can be measured without confounding from stochastic optimization. We use the scaled combinatorial Laplacian
\[
\bar L = L/\|L\|_2
\]
as the graph operator, and apply the same dense-graph scaling to the sparsified Laplacian. The polynomial filter is
\[
p(\bar L)=I-0.6\bar L+0.15\bar L^2.
\]
The map consists of this polynomial graph filter followed by two deterministic feature-mixing stages.

For the training-dynamics experiments of Figure~\ref{fig:training_dynamics_main}, we use the same polynomial-Laplacian graph shift. Dense and sparse trajectories differ only through replacing \(p(\bar L)\) by \(p(\widetilde{\bar L})\), with matched initialization.



For the training-dynamics plots, the number of training epochs is $100$. The learning rates are $0.01$ for the one-layer model and $0.003$ for the two-layer model, with weight decay $10^{-3}$ and gradient clipping.




\paragraph{Teacher models.}
For the synthetic practical experiments, the teacher is a deterministic two-layer GCN trained on the dense graph. The hidden dimension is $48$ for the SBM and $32$ for the geometric graph. Teachers are trained for $120$ epochs with Adam, learning rate $0.01$, and weight decay $5\times 10^{-4}$.

For the real-data experiments, we use a two-layer GCN with ReLU activation and no dropout. When the input feature dimension is at most $64$, we use hidden dimension $48$, which is the case for the datasets considered here. In the full runs, teachers are trained for $20$ epochs with Adam, learning rate $0.01$, and weight decay $5\times 10^{-4}$.

\subsection{Sparsification details}
\label{app:sparsifier}

\paragraph{Synthetic graphs.}
In the controlled synthetic experiments, we use exact effective-resistance sampling on the weighted SBM and geometric graphs. The sampling budget is the number of edge samples drawn with replacement before reweighting. We parameterize it as
\[
q=\lfloor c\,n\log n\rfloor,
\]
where larger $c$ draws more samples and typically yields a denser, lower-distortion sparsifier.
Rather than reporting only a fixed budget sweep, we search over a grid of budget multipliers $c$ and, for each prescribed target level, select a sparsified graph whose realized empirical distortion $\epsilon_{\mathrm{emp}}$ is closest to that target. Each selected target level is then repeated with multiple sparsifier draws. In the base setting we use $5$ repetitions per selected target level.

For each sparsified graph $\widetilde G$, we compute the empirical graph distortion $\epsilon_{\mathrm{emp}}$ from the generalized eigenvalue envelope of the dense and sparse Laplacians on the orthogonal complement of the all-ones vector. Specifically, if
\[
\lambda_{\min} \le \frac{x^\top \widetilde L x}{x^\top L x} \le \lambda_{\max}
\qquad \text{for all } x \perp \mathbf{1},
\]
then we report
\[
\epsilon_{\mathrm{emp}} := \max\{1-\lambda_{\min},\,\lambda_{\max}-1\}.
\]



\paragraph{Real datasets.}
For real-data experiments, we use approximate effective-resistance sampling, where resistance scores are estimated using a truncated Laplacian eigenspace. In the full runs, the approximation rank is $16$, the budget multipliers are
\begin{equation}
c \in \{0.45,0.60,0.80,1.00,1.25,1.50\},
\end{equation}
and each budget is repeated $3$ times with different sparsifier seeds. The empirical graph distortion is approximated by probing quadratic-form distortion with random test vectors.

\subsection{Metric definitions}
\label{app:metrics}

On synthetic graphs, we report filter error, relative hidden representation error, and relative hidden Gram error. On real data, the primary geometry metric is relative hidden Gram distortion,
\begin{equation}
D_{\mathrm{Gram}}(Z,\widetilde Z)
=
\frac{\|ZZ^\top-\widetilde Z\widetilde Z^\top\|_F}
{\|ZZ^\top\|_F}.
\end{equation}

To measure local geometry preservation, we compute mean $20$-nearest-neighbor overlap on held-out nodes. 

Let $z_i^\top$ denote the $i$th row of $Z$, i.e., the hidden embedding of node $i$. We define the similarity between nodes $i$ and $j$ by the cosine similarity of their hidden embeddings,
\[
\mathrm{cos}(z_i,z_j) := \frac{\langle z_i,z_j\rangle}{\|z_i\|_2\|z_j\|_2}.
\]
Let $N_k(i;Z)$ denote the set of the $k$ nodes with largest cosine similarity to node $i$ in the embedding $Z$. The nearest neighbor overlap is given by
%Let $N_k(i;Z)$ denote the $k$-nearest neighbors of node $i$ in hidden space under cosine similarity. \textcolor{red}{[What do you mean here? like cosine similarity of the entries of $Z$ associated with each node or something? maybe a quick definition?]}Then
\begin{equation}
D_{\mathrm{NN}@k}(Z,\widetilde Z)
=
\frac{1}{|S|}
\sum_{i\in S}
\frac{|N_k(i;Z)\cap N_k(i;\widetilde Z)|}{k},
\end{equation}
with $k=20$ in all reported experiments. In the real-data script, these geometry metrics are evaluated on a random subsample of at most $500$ held-out test nodes to keep the computation lightweight. %\textcolor{red}{[Same sort of question: what is a ``held-out'' node? is this common term in CS or just network sci?]}

\paragraph{Compute resources.}
All experiments were run on a workstation with an Intel Core i7-10750H CPU, 16\,GB RAM, and an NVIDIA GeForce GTX 1650 Ti with Max-Q Design GPU (4\,GB GDDR6 VRAM). In practice, the provided code ran primarily on CPU, with approximate runtimes of 5–15 minutes per real dataset and about 20–40 minutes for the full synthetic suite. 


\subsection{Additional qualitative and downstream figures}
\label{app:extra_figures}


Figure~\ref{fig:procrustes_app} provides a qualitative view of the dense and sparse hidden embeddings after projection to two principal components and alignment by an optimal orthogonal Procrustes transform. The substantial overlap after alignment indicates that much of the dense--sparse discrepancy is explained by an approximately rigid transformation rather than severe geometric collapse.

\begin{figure}[htbp!]
    \centering
    \includegraphics[width=0.7\linewidth]{Figures/fig2_final_procrustes.pdf}
    \caption{\textbf{Procrustes-aligned dense and sparse embeddings.}
    Dense and sparse hidden embeddings on held-out nodes after projection to two principal components and optimal orthogonal alignment of the sparse embedding to the dense one. The substantial overlap after alignment indicates that much of the dense--sparse discrepancy is explained by an approximately rigid transformation rather than severe deformation of the learned representation geometry.}
    \label{fig:procrustes_app}
\end{figure}

\subsection{Knowledge distillation under graph distortion}
\label{app:teacher_student}
This experiment provides a downstream diagnostic complementary to the representation-geometry metrics in the main text.
Here, we study teacher--student transfer on the synthetic graphs. A teacher is first trained on the dense graph. We then sparsify the graph and train a student on the sparse graph using a mixture of cross-entropy loss and KL-based knowledge distillation from the dense teacher. Figure~\ref{fig:teacher_student_app} shows that larger empirical graph distortion leads to larger KL divergence in teacher behavior, and that larger teacher KL is associated with worse KD student accuracy. This provides a downstream interpretation of representation distortion: graph perturbations alter the soft predictive information available for transfer.

\begin{figure}[htbp!]
    \centering
    \includegraphics[width=\linewidth]{Figures/fig4_practical_teacher_student.pdf}
    \caption{\textbf{Teacher--student consequences of graph distortion.}
    On synthetic SBM and geometric graphs, larger empirical graph distortion produces larger teacher KL divergence, and larger KL is associated with lower KD student accuracy. This shows that graph distortion affects not only hidden geometry but also the transferred soft decision structure. Error bars denote one standard deviation across four repeated runs.}
    \label{fig:teacher_student_app}
\end{figure}



\end{document}
